\section{Тестування та результати}

\subsection{Показники масштабованості}

Нехай $T_1$ -- час виконання на одному потоці, $T_p$ -- на $p$ потоках. Прискорення та
ефективність визначаються як
\begin{equation}
S(p) = \frac{T_1}{T_p}, \qquad E(p) = \frac{S(p)}{p}.
\end{equation}

Закон Амдала [10] описує межу прискорення програми, частка $f$ якої виконується послідовно:
\begin{equation}\label{eq:amdahl}
S(p) = \frac{1}{f + (1 - f)/p}, \qquad \lim_{p \to \infty} S(p) = \frac{1}{f}.
\end{equation}
Параметр $f$ для експериментальних даних оцінено методом найменших квадратів за точками
$p \le 12$. Закон Густафсона [12] доповнює цю модель для випадку, коли разом із кількістю
процесорів зростає обсяг задачі; у роботі корпус фіксований, тож застосовується модель
Амдала.

Метрика Карпа -- Флатта [11] -- експериментально визначена послідовна частка:
\begin{equation}\label{eq:kf}
e(p) = \frac{1/S(p) - 1/p}{1 - 1/p}.
\end{equation}
Якщо $e$ не змінюється зі зростанням $p$, обмеження спричинене справді послідовною частиною
програми; якщо $e$ зростає -- накладними витратами паралелізму, що збільшуються разом із
кількістю потоків (синхронізація, конкуренція за спільні ресурси) [11].

\subsection{Методика та середовище}

Усі вимірювання виконано на машині автора в штатному розгортанні Docker Compose
(табл.~\ref{tab:env}). Вимірювання автоматизовано класом \code{Bench}, який запускається
всередині контейнера сервера (\code{bench/run\_all.sh}); сирі результати записуються у
CSV-файли каталогу \code{bench/results}, графіки будує \code{bench/plot.py}.

\begin{table}[H]
\caption{Середовище вимірювань}\label{tab:env}
\centering\begin{tabularx}{\textwidth}{lX}
\toprule
Компонент & Характеристика\\
\midrule
Процесор & AMD Ryzen 5 7500F: 6 ядер, 12 апаратних потоків (SMT), базова частота 3,7~ГГц,
L3 32~МБ\\
Пам'ять, доступна Docker & 15,7~ГіБ (віртуальна машина WSL2)\\
Платформа & Windows, Docker Desktop (WSL2), Docker Engine 28.3.3\\
JVM & Eclipse Temurin 21 (JRE), G1 GC, \code{-XX:MaxRAMPercentage=75}\\
LLM & Ollama, модель qwen2.5:1.5b, \code{OLLAMA\_NUM\_PARALLEL=2}, обчислення на CPU\\
Корпус & 16\,793 документи, 1,62~ГБ, том Docker на файловій системі ext4\\
\bottomrule
\end{tabularx}
\end{table}

Методика вимірювання часу побудови:
\begin{itemize}
\item перед серією виконується прогрівна побудова (JIT-компіляція, кеш сторінок ОС), її
результат не враховується;
\item кожна конфігурація вимірюється тричі; три прогони йдуть раундами, і в кожному раунді
порядок конфігурацій випадково перемішується (фіксоване зерно), щоб повільні дрейфи -- нагрів
процесора, стан кешу -- розподілялися рівномірно, а не зміщували окрему конфігурацію;
\item після кожної побудови викликається \code{System.gc()}, щоб сміття попередньої побудови
не збиралося під час наступної;
\item час вимірюється на сервері від початку обходу каталогу до завершення фінальної фази;
у таблицях наведено медіану трьох прогонів, на графіках -- також мінімум і максимум.
\end{itemize}

Пропускну здатність вимірює вбудований навантажувальний клієнт із замкненим циклом: кожен із
$c$ потоків-клієнтів тримає власне keep-alive з'єднання і надсилає наступний запит одразу
після отримання відповіді. Запити циклічно беруться з набору 20 запитів різної частотності
(\code{bench/queries.txt}). Після 5~с прогріву рахуються запити протягом 20~с (для LLM --
30~с), латентність фіксується для кожного запиту окремо. Обмеження методики: клієнт працює в
тому самому контейнері, що й сервер, і ділить з ним 12 апаратних потоків, тому виміряна
пропускна здатність є нижньою оцінкою можливостей сервера.

\subsection{Модульне, стресове та наскрізне тестування}

Для тестів написано власний мінімальний раннер (анотація \code{@Test}, перевірки
\code{check}/\code{equal}/\code{throwsType}); тести запускаються командою
\code{./build.sh test} і під час кожного \code{docker build}. Усього 38 тестів у 9 класах
(табл.~\ref{tab:tests}), повний прогін триває близько 10~с.

\begin{longtable}{p{0.27\textwidth}p{0.66\textwidth}}
\caption{Склад тестів}\label{tab:tests}\\
\toprule
Клас & Що перевіряється\\
\midrule\endfirsthead
\multicolumn{2}{l}{\textit{Продовження таблиці~\ref{tab:tests}}}\\
\toprule Клас & Що перевіряється\\ \midrule\endhead
\code{ThreadPoolTest} & результати через \code{Future}; справжня паралельність (4 задачі
чекають одна на одну на \code{CountDownLatch}); передача винятків і живучість потоків;
graceful shutdown виконує всі 1\,000 поставлених задач і відхиляє нові; \code{shutdownNow}
повертає невиконані; 8 потоків-продюсерів ставлять 80\,000 задач -- кожна виконується рівно
один раз\\
\code{TreiberStackTest} & 8 потоків по 50\,000 операцій push/pop: жоден елемент не втрачено і не
вийнято двічі\\
\code{IndexBuilderTest} & індекс, побудований паралельно, збігається з послідовним еталоном
(окремий простий алгоритм) для обох токенізаторів, обох реалізацій, 1--16 потоків і
1/3/64 партицій -- 60 побудов; 15 побудов у режимі найбільшої конкуренції (батч з одного
документа, 1--2 партиції, 12 потоків); незалежність ідентифікаторів документів від кількості
потоків; колізії хешів (\code{"Aa"} і \code{"BB"}), регістр, службові слова; порожній корпус\\
\code{TokenizerTest} & регістр, роздільники, довжини, не-ASCII, кирилиця, дедуплікація\\
\code{SearchTest} & BM25: вища частота і рідкісніший терм ранжуються вище, слова моделі
важать менше, детермінізм за рівних оцінок; сніпет за цілим словом, на межі блоків читання\\
\code{QueryExpanderTest} & розбір коректних, «загорнутих» і обірваних відповідей моделі;
очищення термів; кеш; деградація за недоступної моделі; дедлайн 300~мс за відповіді
моделі через 3~с; 12 одночасних запитів (4 різні) -- не більше 4 викликів моделі і не більше
2 одночасно (mock-сервер Ollama)\\
\code{JsonTest} & вкладені значення, екранування, некоректні входи\\
\code{WebSocketTest} & ключ рукостискання за прикладом RFC~6455\\
\code{IntegrationTest} & наскрізні тести на запущеному сервері: контракт усіх ендпоінтів і коди
помилок; 50 запитів в одному keep-alive з'єднанні; 16 клієнтів по 20 пошуків
\textit{під час перебудови індексу} -- усі 320 відповідей ідентичні; WebSocket: три етапи,
два запити в одному сеансі, фрагментоване повідомлення, ping/pong, етап \code{error}, коди
закриття 1000/1002/1003/1007; пошук за повільної моделі завершується в межах дедлайну\\
\bottomrule
\end{longtable}

У Java немає аналога ThreadSanitizer, тому відсутність станів гонки перевіряється стресовими
тестами з порівнянням з еталоном: будь-яка втрата чи дублювання входжень під час
конкурентного злиття змінила б частоти термів і була б виявлена. Тест із перебудовою під
навантаженням перевіряє коректність публікації знімка: жоден пошук не побачив частково
побудованого індексу.

\subsection{Час побудови індексу залежно від кількості потоків}

Результати для обох реалізацій спільної структури з 64 партиціями наведено в
табл.~\ref{tab:index} і на рис.~\ref{fig:index_scaling.png}, \ref{fig:index_efficiency.png}.

\begin{table}[H]
\caption{Час побудови індексу (64 партиції, медіана трьох прогонів)}\label{tab:index}
\centering\small
\begin{tabular}{c|ccccc|cc}
\toprule
& \multicolumn{5}{c|}{partitioned (неблокуюча)} & \multicolumn{2}{c}{striped (блокувальна)}\\
$p$ & $T_p$, с & min--max, с & $S$ & $E$ & $e$ & $T_p$, с & $S$\\
\midrule
\tabinput{tables/index.tex}
\bottomrule
\end{tabular}
\end{table}

\fig{index_scaling.png}{Час побудови індексу та прискорення залежно від кількості потоків}
\fig{index_efficiency.png}{Ефективність та експериментальна послідовна частка Карпа -- Флатта}

На двох -- шести потоках прискорення близьке до лінійного з невеликими втратами: 1,75 на 2
потоках і 4,02 на 6 (ефективність 0,87 і 0,67). Шість -- кількість фізичних ядер процесора; за
більшої кількості потоків два потоки ділять одне ядро (SMT), і приріст різко сповільнюється:
від 6 до 10 потоків прискорення зростає лише з 4,02 до 4,85, далі не зростає зовсім, а на
16 потоках, коли потоків більше, ніж апаратних потоків процесора, -- трохи падає (4,51) через
витрати на перемикання контексту. Найкращий результат -- 2,09~с на 10 потоках.

Апроксимація законом Амдала дає послідовну частку $f = 0{,}125$, тобто межу прискорення
$1/f = 8$. Проте справді послідовні ділянки програми значно менші: обхід каталогу триває
близько 78~мс, а фінальна фаза -- близько 0,1~с (табл.~\ref{tab:phases}), разом менше 2~\%
від $T_1$. Експериментальна частка Карпа -- Флатта не постійна: вона становить 0,08--0,14 до
12 потоків і зростає до 0,17 на 16. За критерієм підрозділу~3.1 це означає, що масштабування
обмежують накладні витрати, які зростають разом із кількістю потоків, а не послідовний код.
Їхню природу показує розклад за фазами.

Різниця між реалізаціями спільної структури за 64 партицій лежить у межах розкиду прогонів:
на 8 і 16 потоках трохи швидша блокувальна, на 6, 10 і 12 -- неблокуюча. За такої кількості
партицій синхронізація не є вузьким місцем (підрозділ~3.7).

\subsection{Розклад часу за фазами}

\begin{table}[H]
\caption{Медіани фаз побудови (partitioned, 64 партиції)}\label{tab:phases}
\centering\small
\begin{tabular}{cccccccc}
\toprule
& \multicolumn{3}{c}{Стіновий час фаз} & \multicolumn{3}{c}{Сумарний час потоків, с} &
Передач\\
$p$ & list, мс & map+merge, с & finalize, мс & читання & токенізація & злиття & фрагментів\\
\midrule
\tabinput{tables/phases.tex}
\bottomrule
\end{tabular}
\end{table}

\fig[0.85\textwidth]{index_phases_cpu.png}{Сумарний час потоків за фазами задачі}

Сумарний (за всіма потоками) час кожної фази задачі показано в табл.~\ref{tab:phases} і на
рис.~\ref{fig:index_phases_cpu.png}. Якби накладних витрат не було, висота стовпчиків
залишалася б сталою: та сама робота лише розподіляється між більшою кількістю потоків.
Насправді сумарний час токенізації зростає з 6,84~с на одному потоці до 9,49~с на шести і
15,41~с на дванадцяти, тобто в 2,25 раза. Токенізація не має жодної синхронізації, тож це
зростання -- ознака конкуренції за апаратні ресурси: за SMT два потоки ділять обчислювальні
блоки й кеші L1/L2 одного ядра, усі потоки -- спільний кеш L3 і шину пам'яті, а частота
процесора за навантаження всіх ядер нижча, ніж за одного. Злиття зростає з 2,47 до 5,99~с з
тих самих причин плюс через частіші промахи кешу під час вставок у спільні хеш-таблиці.
Читання файлів займає 0,5--1,8~с сумарно і вузьким місцем не є: після перенесення корпусу в том
Docker (підрозділ~2.1) задача стала обчислювальною.

\subsection{Оптимізація токенізації}

Перша версія задачі створювала об'єкт \code{String} і робила пошук у \code{HashMap} для
кожного слова. Розклад за фазами для неї показав ще стрімкіше зростання сумарного часу
токенізації -- з 10,78~с до 26,13~с на 12 потоках (×2,42) і 37,40~с на 16 потоках. Друга
версія (\code{BatchTermTable}, підрозділ~2.7) хешує байти на місці й створює рядок лише для
нових у батчі термів. Порівняння виконано в тій самій сесії тим самим інструментом
(табл.~\ref{tab:mapper}, рис.~\ref{fig:mapper_compare.png}).

\begin{table}[H]
\caption{Порівняння версій токенізації (partitioned, 64 партиції)}\label{tab:mapper}
\centering\begin{tabular}{cccccc}
\toprule
& \multicolumn{2}{c}{v1: String на слово} & \multicolumn{2}{c}{v2: байтова таблиця} & \\
$p$ & $T_p$, с & токенізація, с & $T_p$, с & токенізація, с & $T^{v1}_p / T^{v2}_p$\\
\midrule
\tabinput{tables/mapper.tex}
\bottomrule
\end{tabular}
\end{table}

\fig{mapper_compare.png}{Час побудови і прискорення двох версій токенізації}

Нова версія швидша в 1,28--1,48 раза за будь-якої кількості потоків: $T_1$ скоротився з
14,28 до 10,12~с, а найкращий час -- з 2,95 до 2,09~с. Сумарний час токенізації зменшився
приблизно вдвічі (на 12 потоках -- з 26,13 до 15,41~с). Прискорення відносно власного $T_1$ в
обох версій майже однакове (рис.~\ref{fig:mapper_compare.png}, праворуч), тобто оптимізація
зменшила обсяг роботи на кожному потоці, але не усунула апаратну конкуренцію, описану в
підрозділі~3.5. Цей результат ілюструє загальне правило: зменшення обсягу виділення пам'яті в
гарячому циклі дає виграш, який множиться на всі потоки, тоді як заміна механізму
синхронізації за малої конкуренції не дає майже нічого (підрозділ~3.7).

\subsection{Вплив кількості партицій і механізму синхронізації}

Щоб навмисно створити конкуренцію за спільну структуру, побудову повторено з 1, 8 і 256
партиціями (смугами) для обох реалізацій (табл.~\ref{tab:partitions},
рис.~\ref{fig:index_partitions.png}, \ref{fig:lock_wait.png}).

\begin{table}[H]
\caption{Час побудови (с) залежно від кількості партицій; синхронізація на 12 потоках}
\label{tab:partitions}
\centering\small
\begin{tabular}{llcccccl}
\toprule
& & \multicolumn{4}{c}{$T_p$ за $p$ потоків, с} & Злиття & Синхронізація\\
Реалізація & $P$ & 1 & 4 & 8 & 12 & (12), с & (12 потоків)\\
\midrule
\tabinput{tables/partitions.tex}
\bottomrule
\end{tabular}
\end{table}

\fig{index_partitions.png}{Час побудови за різної кількості партицій}
\fig[0.8\textwidth]{lock_wait.png}{Сумарний час очікування блокувань (striped)}

З однією партицією обидві реалізації перестають масштабуватися вже після 4 потоків:
3,4--3,7~с на 8 і 12 потоках проти 2,1~с з 8 і більше партиціями. Але причини різні.
Блокувальна реалізація з одним замком на 12 потоках сумарно чекає на нього 22,4~с -- потоки
простоюють, а сумарний час фази злиття зростає до 26,1~с. Неблокуюча реалізація не чекає
зовсім: зайнятість партиції фіксувалася 477 разів, і щоразу фрагмент передавався власнику, а
потік ішов далі. Проте вся робота злиття все одно виконується послідовно -- тим потоком, що
володіє єдиною партицією, -- тож стіновий час той самий. Неблокуючий алгоритм гарантує
прогрес і відсутність простою, але не робить послідовну за природою операцію паралельною: це
наочна ілюстрація закону Амдала, де послідовною частиною є злиття в одну структуру.

Вже 8 партицій прибирають цей ефект для обох реалізацій: очікування на замках падає до
1,5~с (на 12 потоків за 2,1~с роботи), час побудови -- до 2,1~с. Подальше збільшення
кількості партицій до 64 і 256 не прискорює побудову, а в неблокуючій реалізації збільшує
кількість фрагментів і передач (до 6\,330 на 256 партицій), що лише трохи додає накладних
витрат. Отже, за достатньої кількості партицій механізм синхронізації спільної структури
перестає впливати на результат, і масштабування визначається обчислювальною частиною задачі
(підрозділ~3.5). За замовчуванням використовується неблокуюча реалізація з 64 партиціями:
вона ніколи не блокує потік, а її продуктивність не гірша за блокувальну.

\subsection{Пропускна здатність за паралельних клієнтів}

\begin{table}[H]
\caption{Пропускна здатність пошуку без розширення (20~с після 5~с прогріву)}
\label{tab:throughput}
\centering\begin{tabular}{cccccc}
\toprule
Клієнтів & Запитів/с & p50, мс & p95, мс & p99, мс & Помилок\\
\midrule
\tabinput{tables/throughput.tex}
\bottomrule
\end{tabular}
\end{table}

\fig{throughput.png}{Пропускна здатність і латентність пошуку залежно від кількості клієнтів}

Пропускна здатність (табл.~\ref{tab:throughput}, рис.~\ref{fig:throughput.png}) зростає
майже лінійно до 8 клієнтів: 1\,247 запитів/с на одному, 9\,172 на восьми (×7,4) за майже
незмінної медіанної латентності 0,76--0,81~мс. Це очікувано: пошукові потоки лише читають
незмінний знімок індексу, не мають спільних записуваних даних і не синхронізуються між собою.
Максимум -- 12\,004 запитів/с -- досягається на 16 клієнтах; далі пропускна здатність
стабілізується на рівні 11\,500, а латентність зростає (p99 -- 13,1~мс на 32 і 31,7~мс на 64
клієнтах): усі 12 апаратних потоків зайняті, і нові клієнти лише довше чекають на процесор.
Частину процесорного часу при цьому споживає сам навантажувальний клієнт, що працює на тій
самій машині. Помилок не зафіксовано.

\subsection{Вузьке місце LLM-розширення}

\begin{table}[H]
\caption{Пошук з розширенням через LLM}\label{tab:llm}
\centering\begin{tabular}{lcccc}
\toprule
Режим & Клієнтів & Запитів/с & p50, мс & p95, мс\\
\midrule
\tabinput{tables/llm.tex}
\bottomrule
\end{tabular}
\end{table}

\fig{throughput_llm.png}{Пропускна здатність і латентність пошуку з розширенням (без кешу)}

Виклик моделі на CPU триває близько 0,6~с, що на три порядки довше за сам пошук
(табл.~\ref{tab:llm}, рис.~\ref{fig:throughput_llm.png}). Без кешу пропускна здатність
становить 1,6~запиту/с для одного клієнта і 2,3--2,4 запиту/с для двох і більше: семафор
пропускає до моделі не більше двох запитів одночасно, а Ollama обробляє не більше двох
паралельно. Подальше збільшення кількості клієнтів лише подовжує чергу, і латентність зростає
пропорційно: 0,89~с для двох, 1,69~с для чотирьох, 3,32~с для восьми клієнтів. Це узгоджується
із законом Літтла $L = \lambda W$: за $\lambda \approx 2{,}4$~запиту/с і $L = 8$ запитів у
системі $W \approx 3{,}3$~с. Модель -- типовий послідовний ресурс конвеєра, і вирішальним є
не те, як паралельно працює решта системи, а те, щоб очікування на модель не займало
спільні ресурси. Тому в реалізації (підрозділ~2.11) очікування обмежене дедлайном, знімок
індексу береться лише після розширення, а однакові запити об'єднуються.

Кешування розширень повністю прибирає модель із шляху повторних запитів: з кешем пропускна
здатність становить 959 запитів/с на одному клієнті і 8\,634 на 32 клієнтах, що близько до
пошуку без розширення (решту часу займають пошук за розширеним набором термів і кеш під
замком).

\subsection{Поведінка пам'яті}

\begin{table}[H]
\caption{Пам'ять під час шести послідовних перебудов на 12 потоках, МіБ}\label{tab:memory}
\centering\begin{tabular}{c|ccc|ccc}
\toprule
& \multicolumn{3}{c|}{без \code{System.gc()}} & \multicolumn{3}{c}{з \code{System.gc()}}\\
Перебудова & heap used & committed & RSS & heap used & committed & RSS\\
\midrule
\tabinput{tables/memory.tex}
\bottomrule
\end{tabular}
\end{table}

\fig{memory.png}{Пам'ять JVM під час повторних перебудов індексу}

Готовий індекс займає в купі близько 350~МіБ (табл.~\ref{tab:memory}, рис.~\ref{fig:memory.png}):
33,2~млн упакованих входжень по 8~байт (253~МіБ) плюс словник. Без явного збирання сміття
після перебудови зайнята купа становить 1,7--3,2~ГіБ: G1 не поспішає збирати старий індекс і
тимчасові структури побудови, доки є вільна пам'ять. З викликом \code{System.gc()} після підміни
зайнята купа стабільно повертається до 352~МіБ, тобто витоку пам'яті між перебудовами немає.
RSS процесу при цьому коливається від 2,5 до 5,2~ГіБ: JVM утримує раз виділені сторінки купи
(\code{committed} -- 1,3~ГіБ після збирання, але пікові сторінки під час побудови повертаються
операційній системі лише частково). Звідси практичний висновок для розгортання: на час
перебудови потрібно передбачити пам'ять щонайменше для двох індексів і проміжних таблиць
задач, тобто кілька гігабайтів, хоча сам індекс займає лише 350~МіБ.

\subsection{Функціональні тестові запити}

Для перевірки якості пошуку виконано набір запитів, очікуваний результат яких відомий
наперед (табл.~\ref{tab:queries}; ідентифікатор документа -- номер книжки Project Gutenberg і
номер частини). Без розширення пошук займає менше 1~мс, з розширенням -- 0,55--0,82~с.

\begin{longtable}{p{0.18\textwidth}p{0.25\textwidth}p{0.30\textwidth}p{0.14\textwidth}}
\caption{Результати пошуку на тестових запитах}\label{tab:queries}\\
\toprule
Запит & Очікуваний результат & Фактичний результат (топ-3) & Оцінка\\
\midrule\endfirsthead
\multicolumn{4}{l}{\textit{Продовження таблиці~\ref{tab:queries}}}\\
\toprule Запит & Очікуваний & Фактичний (топ-3) & Оцінка\\ \midrule\endhead
moby dick white whale & книжки «Moby Dick» (PG 15, 2489, 2701) & 15\_0000, 2489\_0003,
2701\_0003 -- три видання «Moby Dick» & відповідає\\
sherlock holmes baker street & оповідання про Шерлока Холмса & 2038\_0001 (збірка з «A Scandal
in Bohemia»), 221\_0000 і 108\_0000 («The Return of Sherlock Holmes») & відповідає\\
tarzan apes jungle & книжки про Тарзана (PG 78, 106) & 78\_0001 «Tarzan of the Apes»,
106\_0002, 106\_0003 «Jungle Tales of Tarzan» & відповідає\\
romeo juliet & п'єса «Romeo and Juliet» (PG 1513) & 1513\_0000, 1513\_0001, 1286\_0004 &
відповідає\\
pride prejudice darcy & «Pride and Prejudice» (PG 1342) & 1342\_0006, 1342\_0003, 1342\_0000 &
відповідає\\
abraham lincoln & біографії Лінкольна & 1815\_0000 «The Boys' Life of Abraham Lincoln»,
1713\_0006, 2517\_0006 & відповідає\\
electricity magnetism & праці з електромагнетизму & 1225\_0001, 1225\_0000 «Faraday As A
Discoverer», 1707\_0004 & відповідає\\
treasure island & «Treasure Island» (PG 120) & 973\_0000 «Howard Pyle's Book of Pirates»,
3029\_0000, 2021\_0009 & PG 120 немає в корпусі; знайдено тематично близьке\\
quixotic & рідкісне слово, точні збіги & 2158\_0012 «The Prime Minister», 1803\_0003,
2158\_0011 & відповідає\\
xyzzyplugh & порожня видача & без розширення -- порожня; з розширенням модель вигадала
«xylophone», «clavier», «melody» і дала нерелевантні документи & див. текст\\
the and of & порожня видача (лише службові слова) & порожня, звернення до моделі не
виконується & відповідає\\
\bottomrule
\end{longtable}

З розширенням видача змінюється змістовно: для «moby dick white whale» модель додала
«captain», «ahab», для «abraham lincoln» -- «president», «union», «republican», і на першому
місці опинився том праць самого Лінкольна (PG 2653); для «treasure island» -- «pirate»,
«captain», «gold», і у видачу потрапив «Peter Pan» (PG 16). Слабке місце розширення
показує безглуздий запит «xyzzyplugh»: модель не відмовляється від відповіді, а вигадує
асоціації за схожістю написання, і запит без жодного збігу отримує непорожню видачу.
Зменшена вага слів моделі (0,4) не допомагає, коли слів користувача в корпусі немає зовсім;
можливе покращення -- не розширювати запит, якщо жоден його терм не знайдено в індексі.
